% ============================================================================
% Minimal example — shows how to cherry-pick LIRN components
% ============================================================================
% You don't have to use everything. Mix and match.
% ============================================================================

\documentclass[report]{lirn}

\preparedfor{Any Client}
\docdate{13 April 2026}

\begin{document}

% --- Just the header --------------------------------------------------------
\lirnheader

% --- Just a title -----------------------------------------------------------
\lirntitle{Quick Technical Note}

% --- Just the meta line -----------------------------------------------------
\lirnmeta

% --- Write anything ---------------------------------------------------------
\section{Observation}

This is a minimal document using LIRN branding. You can inject any component
independently:

\begin{itemize}
  \item \texttt{\textbackslash lirnheader} --- the branded banner
  \item \texttt{\textbackslash lirntitle\{...\}} --- styled document title
  \item \texttt{\textbackslash lirnmeta} --- Prepared for / Date line
  \item \texttt{\textbackslash lirnsubject} --- Subject line
  \item \texttt{\textbackslash lirnprocess} --- Process line
  \item \texttt{\textbackslash lirnrefnote\{Label\}\{Content\}} --- callout box
  \item \texttt{\textbackslash lirnrule} --- horizontal divider
  \item \texttt{\textbackslash lirnend[Report]} --- end-of-document marker
  \item \texttt{\textbackslash lirndisclaimers} --- default disclaimers
  \item \texttt{\textbackslash lirndisclaimers[...]} --- custom disclaimer items
\end{itemize}

% --- A callout anywhere -----------------------------------------------------
\lirnrefnote{NOTE}{%
  You can drop a callout box anywhere in any document.
  Just call \texttt{\textbackslash lirnrefnote} with a label and content.%
}

% --- End marker + disclaimers anywhere --------------------------------------
\lirnend[Note]
\lirndisclaimers

\end{document}
